% !TEX root = ../main.tex
\onehalfspacing
{\fontsize{16pt}{30pt}\selectfont \noindent\textbf{摘要}}
\vspace{0.7cm}
{\fontsize{12pt}{18pt}\selectfont
% replace the paragraph below with your abstract

随着大语言模型智能体广泛通过调用外部工具执行任务，现有 Linux 操作系统缺乏对该过程的原生管控，导致控制平面的核心逻辑完全局限于用户态的网关守护进程中。
这种架构使得本应由内核保障的安全与可靠性属性，过度依赖于单个用户态进程的正确性与存活性。

针对上述痛点，本文设计并实现了 \texttt{linux-mcp}，一种面向 Linux 的内核辅助型大语言模型智能体工具调用控制平面。
该架构将语义解析与工具执行保留在用户态，而将核心控制平面状态（包含已注册的工具与智能体、生效的授权凭证及会话上下文）下沉至内核模块维护。
该机制无需修改现有工具，仅需通过声明式 JSON 清单与对应适配层即可接入。借此设计，授权状态不仅能在守护进程崩溃后实现持久化，还能有效抵御非特权用户态进程的恶意篡改。

本文从性能、抗攻击与故障恢复三个维度，将 \texttt{linux-mcp} 与纯用户态基线及 seccomp 加固版本进行了对比评估。
在针对性控制平面攻击测试集中，\texttt{linux-mcp} 拦截了全部攻击，纯用户态基线则有 84.6\% 的攻击成功绕过，且常见负载下端到端延迟仍在测量噪声范围之内。
扩展攻击实验另发现原内核模块中漏掉了一项 peer-credential 校验，已在不引入可观测延迟的前提下完成修复。
原型代码已开源于 \url{https://github.com/xiheng724/linux_mcp}。

}
\vspace{1.0cm}
{\fontsize{16pt}{30pt}\selectfont \noindent\textbf{关键词}}
\vspace{0.7cm}
{\fontsize{12pt}{18pt}\selectfont
% replace the paragraph below with your keywords

AI智能体，操作系统，工具调用，系统安全

}
